

\emph{Convince your enemy that he will gain very little by attacking you, which will diminish their enthusiasm.} \\ 
     -- Sun Tzu, \emph{The Art of War} \citeme

\section{Introduction}
Since the dawn of time there have been wars, and for just as long, people have tried to figure out how to avoid them. Deterring a potential enemy from attacking requires them to believe that the costs of such an attack outweighs their perceived benefits. In this simplified manner there are two strategies, reduce the potential benefits or increase the costs. The former is often counter-productive, though has been successful for example with the burn land tactic during the second world war, \citeme. The latter, however, requires that the attacker believes the consequences to be large, for example by believing that large scale retaliation would occur -- so-called second strike capabilities. 

These ideas form the baseline for \emph{deterrence} as a military strategy: avoiding an attack by having sufficiently powerful means of guaranteed retaliation \citeme. In this paper we study this deterrence mechanism for \emph{coalitions} as a classification problem using ROC curve analysis. 

In essence, deterrence relies on a push-pull between two fundamental objectives that are, in effect, the strategic analogues of true and false positive classifications \citeme: 
\begin{enumerate}
    \item the resolve to use retaliatory means have to be sufficiently credible, and 
    \item one should limit the possibility of escalation spirals and accidental use. 
\end{enumerate} 
The dynamic between these two mechanisms is especially potent for nuclear deterrence, where the consequences are potentially of an existential nature. For a classical outlook on these fundamental issues see \citeme and \citeme. 
% see, inter alia, Schelling and subsequent deterrence theory.

NATO is the only proper nuclear alliance -- in the sense that of being a multinational alliance and not a bilateral agreements -- where the deterrence capabilities of individual nations are also shared to allies. However, NATO does not own nuclear weapons as a collective entity, nor does it operate a jointly owned nuclear command structure with shared launch authority. For example, U.S. weapons stationed in Europe remain under exclusive U.S. custody and control, and any nuclear mission requires authorization from the relevant national leaders in addition to political approval through NATO's Nuclear Planning Group (NPG). \citeme

Hence, there is currently no coalition of states in the world that jointly own both the nuclear deterrence capabilities and the command structure to act upon threats, as the historical proposals for a NATO multinational nuclear force never materialized \citeme. However, such coalitions might become more politically viable now that arms-control frameworks, such as the New START treaty \citeme, are intensifying the European discussions about burden sharing and the possible future role of France or the UK in European nuclear deterrence. 

Currently, of course, the Non-Proliferation Treaty \citeme serves as a block to the creation of such nuclear alliances. However, there is room to break the treaty in the case of existential threats, or, alternatively, one could consider other deterrence technologies not limited by such restrictions. In the advent of such coalition then, the additional problem of \emph{collective} rather than \emph{unilateral} decision-making arises to complicate the above two problems. The problem of \emph{credible resolve} becomes a problem of how the coalition's states aggregate their individual resolves into a joint decision, and the problem of \emph{escalation control} becomes a problem of not allowing this joint decision to be too easily affected by individual potential errors. Hence this joint institutional decision-structure introduces a trade-off between the two fundamental problems. 

We can then phrase the motivating question for the present paper as: How does the institutional design of a coalition affect the deterrence mechanism and its effectiveness? 

Such aggregations of individual choices is in social choice theory measured by a \emph{social choice function}, see \citeme. To answer this question then, we develop a simple Bayesian decision model inspired by signalling games. This model is dependent on weighted threshold choice function $f$ that aggregates the coalitions votes on retaliation in to a shared resolve. 
 
To understand the interactions between deterrence and institutional design, we compute simple attack conditions, and use this to formulate the choice of institutional design for the coalition as a \emph{binary classification problem}. Each institutional design has an associated empirical Receiver Operating Characteristic (ROC) curve -- see \citeme -- for each information environment, formalized by vectors of probabilities of resolves and false positives for the coalition members. By studying the statistics of these ROC curves for a small hypothetical coalition of four members, we can study which designs that perform well and which do not. This does not capture the full strategic payoffs for the deterrence coalition, but gives a tractable and computable way to compare social choice functions on their detection abilities. 

A voting scheme that performs well on a variety of probability distributions for retaliation and false alarms is, perhaps not surprisingly, the unbiased one, where each coalition members has equal say. Close behind follows a scheme where each coalition members is given a weighted vote according to some relevant capabilities, like ISR. Across our stylized information environments, the worst performer is the dictatorial scheme, where only one country has voting power over the use of deterrence technology.

Finally we use Youden's $J$-statistic to study which thresholds that give the best tradeoff between credible deterrence and escalation control, concluding that the best threshold usually is $2$ for our small hypothetical coalition of four members. 

Our goal with this paper is to provide a first step toward understanding the interaction between institutional design and collective deterrence by introducing, studying and exploring a novel binary classification-based framework. We deliberately adopt a simplified game model to isolate the informational tradeoffs inherent in collective decision-making -- extensions incorporating more strategic behavior, correlated information, and proofs of optimality will be studied in future work. 

\subsection{Contributions}
This paper has three main contributions towards the understanding of institutional design in collective deterrence: 
\begin{enumerate}
    \item introducing a novel classification-based framework for analyzing collective deterrence, with an easily understandable visual element;
    \item providing a tractable model for linking institutional design to deterrence credibility and escalation risk;
    \item illustrating the framework through simulations of small coalitions under varying informational environments.
\end{enumerate}

Our results should be interpreted as illustrative patterns and not general results, as they are based on statistical analysis for a small coalition. The results, however, give clear directions for future work on optimal institutional designs using a more general and proof-based approach. 
